Software documentation, which provides detailed and clear descriptions of source code, is crucial for repository comprehension. Researchers have developed various automated methods, as manual documentation writing is labor-intensive. With the advancement of large language models (LLMs), these methods extend from isolated code snippets to the entire repository, leveraging global semantic context for comprehensive summaries. However, existing benchmarks for evaluating software documentation suffer from two fundamental limitations. (1) They lack repository-level analysis, assessing in a fragmented manner that overlooks overall documentation quality. (2) They depend on unreliable evaluation strategies. While LLM-as-a-judge methods are widely adopted, their reliability is compromised by vaguely defined evaluation criteria and limited repository-level knowledge.

To address these limitations, we propose a novel benchmark for evaluating repository-level \textbf{S}oft\textbf{W}are \textbf{D}ocumentation, named \textbf{\benchmark}. Our evaluation strategy is inspired by documentation-driven development, where higher-quality documentation enables more effective repository comprehension. Based on this strategy, we propose to regard LLMs as repository developers and evaluate the documentation quality through the process of the LLMs' understanding and implementing functionalities, instead of directly prompting LLMs for evaluation. To ensure the reliability of evaluation results, documentation quality is assessed by functionality-driven question answering (QA) tasks.
Specifically, \benchmark introduces three interconnected QA tasks: (1) \textbf{Functionality Detection}, aiming to assess whether a specific functionality exists in the documentation.
(2) \textbf{Functionality Localization}, aiming to evaluate the capability in accurately locating functionality-related files. (3) \textbf{Functionality Completion}, aiming to measure the comprehensiveness of implementation details of the functionalities.
The construction pipeline for \benchmark involves three stages: we first mine high-quality Pull Requests (PRs) through multi-step filtering, then enrich them with diverse repository-level context, and finally leverage this rich context to meticulously craft the tasks. This rigorous process yields the final benchmark of 4,170 entries across three QA tasks. Extensive experiments reveal that there still exist limitations in current repository-level documentation generation methods, and highlight that source code offers complementary value to software documentation. Besides, documentation generated by the best-performed method improves the issue-solving rate of SWE-Agent, one popular issue-fixing approach, by 20.00\%, highlighting the practical value of high-quality documentation in supporting documentation-driven development.